\documentclass{article}
\usepackage{amsmath, amssymb, amsthm}
\usepackage[margin=2.5cm]{geometry}
\usepackage{xr}
% Resolve \ref/\eqref to labels defined in the main notes (compile notes.tex
% first so that notes.aux exists).
\externaldocument{notes}

\theoremstyle{definition}
\newtheorem*{remark}{Remark}

\title{Onset-anchored models with delayed transmission}
\author{}
\date{}

\begin{document}
\maketitle

These notes are a companion to the main methodological notes
(\texttt{notes.tex}) and develop a fully \emph{onset-anchored} branching-process
model, in which transmission is referenced to the infector's own symptom onset
rather than to their (unobserved) infection time.  We build on the main notes
throughout: the SSE, SSI and Poisson models of \S\ref{sec:models}, the
incubation period and infection-to-onset mapping of \S\ref{sec:bridge}, the
probability of a major outbreak under model uncertainty of
\S\ref{sec:pmo_uncertain}, and the real-time particle filter of
\S\ref{sec:realtime}.  Cross-references below (\S and equation numbers) are to
the main notes unless stated otherwise.

% ===========================================================================
\section{Onset-anchored SSE and SSI models}
\label{sec:onset_models}
% ===========================================================================

The bridge model of \S\ref{sec:bridge} keeps the generation-time renewal of
\S\ref{sec:models} and adds the incubation period only as an observation delay.
Here we go further and anchor transmission directly to the infector's own
\emph{symptom onset}.  This is appropriate for a disease with negligible
pre-symptomatic transmission --- Ebola virus disease is our motivating example
--- for which the onset is the natural reference point for a case's
infectiousness.

As in \S\ref{sec:bridge} let $D_t \geq 0$ be the number of symptom onsets on
day~$t$ (the observed quantity) and $J_t$ the number of new infections on
day~$t$, with the incubation weights $\iota_a$ ($a \geq 1$) mapping infections
forward to onsets through the multinomial split \eqref{eq:bridge_onsets}.  In
place of the generation interval we introduce the \emph{time from symptom onset
to transmission} (TOST) weights $\tau_s$, $s = 0, 1, 2, \ldots$ with
$\sum_{s \geq 0} \tau_s = 1$, giving the probability that a transmission event
occurs $s$ days after the infector's own onset.  Pre-symptomatic transmission
would place mass on $s < 0$ and is assumed negligible.  The force of infection is
now the TOST-weighted sum of past \emph{onset} cohorts.  In the onset-anchored
SSE model,
\begin{equation}\label{eq:onset_sse}
  \lambda_t = \sum_{s=0}^{\infty} \tau_s\, D_{t-s},
  \qquad
  J_t \mid \lambda_t \;\sim\;
  \mathrm{NB}\!\left(\text{mean} = R_0 \lambda_t,\; \text{disp} = k\lambda_t
  \right),
\end{equation}
while in the onset-anchored SSI model each onset cohort carries an aggregate
latent infectivity and infections arise as a Poisson process driven by past
infectivity,
\begin{equation}\label{eq:onset_ssi}
  Y_t \mid D_t \;\sim\; \mathrm{Gamma}\!\left(\text{shape} = k D_t,\;
  \text{rate} = k\right),
  \qquad
  \lambda_t = \sum_{s=0}^{\infty} \tau_s\, Y_{t-s},
  \qquad
  J_t \mid \lambda_t \;\sim\; \mathrm{Poisson}\!\left(R_0 \lambda_t\right),
\end{equation}
exactly as in \S\ref{sec:ssi} but with the serial-interval weights $w_s$ replaced
by the TOST weights $\tau_s$ and the driving cohorts taken to be onsets.  As
$k \to \infty$ both reduce to a common onset-anchored Poisson model with
$J_t \sim \mathrm{Poisson}(R_0 \lambda_t)$ and
$\lambda_t = \sum_{s \geq 0} \tau_s D_{t-s}$.  The generative process alternates
$D_t \to \{Y_t\} \to J_t \to D_{>t}$ --- with $J_t$ mapped forward to future
onsets by \eqref{eq:bridge_onsets} --- seeded by a single index case with onset
on day~0 ($D_0 = 1$).

\paragraph{Indexing convention and causal ordering.}
Because the incubation weights are supported on $a \geq 1$, an infection on
day~$t$ cannot produce an onset before day~$t+1$.  Consequently the onset cohort
$D_t$ --- and, in the SSI model, its infectivity $Y_t$ --- is completely
determined by infections on days strictly before $t$, so it is fully realised
before day~$t$'s own transmission is evaluated in
\eqref{eq:onset_sse}--\eqref{eq:onset_ssi}.  This well-orders the recursion.  The
TOST weights, by contrast, are supported on $s \geq 0$, so a case may transmit on
its own onset day.  At most one of the two distributions may place mass at
lag~0: if both did, a chain of same-day infection, onset and transmission could
occur within a single day and the recursion would no longer be well-defined.

\paragraph{Relationship to the serial interval.}
Consider an infector with onset on day~$u$.  It transmits on day $u + S$ with
$S \sim \tau$; the infectee, infected on day $u + S$, has onset on day
$u + S + A'$ with $A' \sim \iota$.  The onset-to-onset serial interval is
therefore $S + A'$, distributed as the convolution $\tau * \iota$ with mean
$\mathbb{E}[\tau] + \mathbb{E}[\iota]$.  Choosing the TOST and incubation means so
that they sum to the serial-interval mean used in \S\ref{sec:models} yields an
onset-anchored model that is roughly comparable to the infection-anchored models,
while making explicit the intermediate onset-to-transmission and
infection-to-onset delays.  (This differs from the bridge model of
\S\ref{sec:bridge}, where the onset-to-onset interval is the generation interval
convolved with the \emph{difference} of two incubation periods; there the
incubation enters only through observation, whereas here the TOST anchors
transmission itself.)

% ===========================================================================
\section{Probability of a major outbreak}
\label{sec:onset_pmo}
% ===========================================================================

Given an observed onset history $D_0, \ldots, D_r$ (with $D_0 \geq 1$), we again
ask for the probability of a major outbreak.  Exactly as in the bridge model
(\S\ref{sec:bridge}), and unlike the infection-anchored SSE model of
\S\ref{sec:qr_sse}, the observed onsets do not determine the state needed to
project the epidemic forward: at the observation horizon there is a pipeline of
infected-but-not-yet-symptomatic individuals, whose number and future onset days
are latent, together with (in the SSI model) the latent infectivities.  We
therefore estimate the probability of a major outbreak by the particle filter of
\S\ref{sec:realtime}, matching the observed onset cohorts $D_1, \ldots, D_r$;
equivalently, for a single retrospective horizon, by the rejection scheme of
\S\ref{sec:simulation} applied to the onset process.  A trajectory is resolved as
a major outbreak when a single-day onset count reaches the threshold, and as
extinct once no infected individuals remain in the incubation pipeline
\emph{and} no onsets have occurred within the TOST support of recent days, so
that no further transmission is possible; tracking the pipeline (infections
generated minus onsets realised) is necessary because a recent run of zero onsets
alone does not preclude onsets already scheduled from earlier infections.

This scheme applies unchanged to the onset-anchored SSE, SSI and Poisson models.
Closed-form expressions for the residual extinction probability, analogous to
those of \S\ref{sec:qr}, are likely tractable --- the SSE and Poisson offspring
still decompose independently, and the SSI infectivities admit the same
conditioning arguments --- but are not pursued here; the simulation-based
estimator suffices for the exploratory comparisons of interest.

\paragraph{Real-time monitoring and model averaging.}
The bootstrap particle filter of \S\ref{sec:realtime} applies directly, advancing
onset trajectories one week at a time and retaining the particles whose newly
realised onset cohort matches the observation.  Because the incubation delay is
at least one time unit, the week-$t$ onset cohort is already determined by the
state carried in each particle, so conditioning on it requires no fresh
simulation.  The mean matching fraction at each week is the sequential Monte
Carlo estimate of the one-step predictive probability
$\mathbb{P}_M(D_w \mid D_{0:w-1})$, and its running product is the marginal
likelihood $L_M$ of \S\ref{sec:pmo_uncertain}; feeding the per-model evidences
through \eqref{eq:pi_post_general}--\eqref{eq:pmo_avg_general} gives the posterior
model probabilities and the model-averaged probability of a major outbreak,
tracked in real time as the outbreak unfolds.

\paragraph{Reporting delays (future extension).}
A natural refinement retains this onset-process model but adds a reporting delay:
at the time of analysis some symptom onsets have not yet been reported, so the
observed history is a right-truncated and thinned version of the true onset
history.  The same simulation-based framework accommodates this by matching only
the reported subset of onsets (or by marginalising over unreported recent
onsets) when accepting trajectories.  We note it here as a direction for future
work and do not develop it further.

\end{document}
